\subsection{Minors, Cofactors, and Rank}

\begin{definition} \label{def:minors-cofactors-and-adjugate} \hlt{Minors, Cofactors, and Adjugate}\\
For $A\in\operatorname{Mat}(n,n,\mathbb{F})$, let $M_{ij}$ delete row $i$ and column $j$ from $A$; set $\det(\emptyset_{0\times0})=1$. Define $C_{ij}\defeq(-1)^{i+j}\det(M_{ij})$.\\
Then $C=[C_{ij}]$ is the \hlt{cofactor matrix}, and $\operatorname{adj}(A)\defeq C^t$ is the \hlt{adjugate}.
\end{definition}

\begin{lemma} \label{lem:cofactor-expansion} \hlt{Cofactor Expansion}\\
For $A=[a_{ij}]\in\operatorname{Mat}(n,n,\mathbb{F})$, $\det A=\sum\limits_{k=1}^{n}a_{ik}C_{ik}$ and $\det A=\sum\limits_{k=1}^{n}a_{kj}C_{kj}$.
\begin{proof}
Using $\det A=\sum\limits_{\sigma\in S_n}\operatorname{sgn}(\sigma)\prod\limits_{r=1}^n a_{r,\sigma(r)}$, group terms by $\sigma(i)=k$.\\
For fixed $k$, the remaining permutation sum equals $(-1)^{i+k}\det(M_{ik})=C_{ik}$, giving the row formula. Applying this to $A^t$ and using Theorem \ref{thm:determinant-basic-laws} gives the column formula.\\
\hlt{Structural Commentary}: Cofactor expansion is the big determinant formula collected around one row or column.\\
\hlt{Key Insight}: Fix one entry; the rest is its minor.
\end{proof}
\end{lemma}

\begin{remark} \label{rem:bordered-determinant-adjugate-form} \hlt{Adjugate Form of the Bordered Formula}\\
Proposition \ref{prop:bordered-determinant-formulas} gives the adjugate form $\det
\begin{bmatrix}
A & \mathbf{x}\\
\mathbf{y}^t & c
\end{bmatrix}
=
c\det(A)-\mathbf{y}^t\operatorname{adj}(A)\mathbf{x}$.
No invertibility of $A$ is required.
\end{remark}

\begin{theorem} \label{thm:adjugate-identity-and-cramers-rule} \hlt{Adjugate Identity and Cramer's Rule}\\
Let $A\in\operatorname{Mat}(n,n,\mathbb{F})$ and let $\mathbf{b}\in\mathbb{F}^n$. Then $A\operatorname{adj}(A)=\operatorname{adj}(A)A=(\det A)I_n$. If $\det A\neq0$, then, for $A\mathbf{x}=\mathbf{b}$, $A^{-1}=\frac{1}{\det A}\operatorname{adj}(A)$ and $x_j=\frac{\det A_j(\mathbf{b})}{\det A}$ for $1\leq j\leq n$, where $A_j(\mathbf{b})$ replaces column $j$ of $A$ by $\mathbf{b}$.
\begin{proof}
Since $(\operatorname{adj}A)_{kj}=C_{jk}$, $(A\operatorname{adj}A)_{ij}=\sum\limits_k a_{ik}C_{jk}$.\\
For $i=j$ this is $\det A$ by Lemma \ref{lem:cofactor-expansion}; for $i\neq j$ it is the cofactor expansion of a matrix with rows $i$ and $j$ equal, hence $0$. Thus $A\operatorname{adj}(A)=(\det A)I_n$. The same argument applied to $A^t$ gives $\operatorname{adj}(A)A=(\det A)I_n$. If $\det A\neq0$, division gives the inverse formula. Finally, $(\operatorname{adj}A\,\mathbf{b})_j=\sum_i C_{ij}b_i=\det A_j(\mathbf{b})$, by expansion down column $j$, so $x=A^{-1}\mathbf{b}$ gives Cramer's rule.\\
\hlt{Structural Commentary}: Cofactors assemble all one-column replacements of $A$ into a single inverse formula.\\
\hlt{Key Insight}: Off-diagonal products vanish because they are determinants with two equal rows or columns.
\end{proof}
\end{theorem}

\begin{theorem} \label{thm:similarity-preserves-trace-and-determinant} \hlt{Trace and Determinant Are Similarity Invariants}\\
If $A,B\in\operatorname{Mat}(n,n,\mathbb{F})$ are similar, then $\operatorname{tr}(A)=\operatorname{tr}(B)$, $\det(A)=\det(B)$.
\begin{proof}
Write $B=PAP^{-1}$ with $P$ invertible.\\
Proposition \ref{prop:trace-cyclicity-two-factors} gives $\operatorname{tr}(B)=\operatorname{tr}\bigl(P(AP^{-1})\bigr) =\operatorname{tr}\bigl((AP^{-1})P\bigr) =\operatorname{tr}(A)$.\\
Theorem \ref{thm:determinant-basic-laws} gives $\det(B)=\det(P)\det(A)\det(P^{-1})=\det(A)$.\\
\hlt{Structural Commentary}: Trace and determinant ignore the coordinate system used to write an operator.\\
\hlt{Key Insight}: Conjugating by $P$ inserts $P$ and $P^{-1}$, which cancel inside trace and determinant.
\end{proof}
\end{theorem}

\begin{principle} \label{princ:trace-determinant-coordinate-free} \hlt{Trace and Determinant Are Coordinate-Free}\\
Trace and determinant are first computed from a matrix, but their real structural meaning is operator-level: changing basis conjugates the matrix, and conjugation preserves both quantities.
\end{principle}

\begin{theorem} \label{thm:operator-matrix-trace-determinant-basis-independent} \hlt{Trace and Determinant of an Operator Are Basis-Independent}\\
Let $\mathcal{V}$ be finite-dimensional, let $T\in\LL(\mathcal{V})$, and let $\beta,\gamma$ be bases of $\mathcal{V}$.\\
Then $\operatorname{tr}\bigl([T]_{\beta\leftarrow\beta}\bigr) =\operatorname{tr}\bigl([T]_{\gamma\leftarrow\gamma}\bigr)$, $\det\bigl([T]_{\beta\leftarrow\beta}\bigr) =\det\bigl([T]_{\gamma\leftarrow\gamma}\bigr)$.
\begin{proof}
By Theorem \ref{thm:change-of-basis-formula}, the two matrices are similar: $[T]_{\gamma\leftarrow\gamma} =C_{\gamma\leftarrow\beta}[T]_{\beta\leftarrow\beta}C_{\beta\leftarrow\gamma}$.\\
Theorem \ref{thm:similarity-preserves-trace-and-determinant} now gives equality of traces and determinants.\\
\hlt{Structural Commentary}: Basis changes conjugate matrices; trace and determinant remain operator invariants.\\
\hlt{Key Insight}: Similar matrices have the same trace and determinant.
\end{proof}
\end{theorem}

\begin{definition} \label{def:operator-trace-and-determinant} \hlt{Trace and Determinant of an Operator}\\
Let $\mathcal{V}$ be finite-dim, let $T\in\LL(\mathcal{V})$. Choose any basis $\beta$ of $\mathcal{V}$. Define $\operatorname{tr}T\defeq\operatorname{tr}\bigl([T]_{\beta\leftarrow\beta}\bigr)$, $\det T\defeq\det\bigl([T]_{\beta\leftarrow\beta}\bigr)$.\\
Theorem \ref{thm:operator-matrix-trace-determinant-basis-independent} shows these definitions do not depend on chosen basis.
\end{definition}

\begin{proposition} \label{prop:isometry-determinant-absolute-one} \hlt{Unitary Matrices and Isometries Have Determinant of Modulus One}\\
Let $U\in\operatorname{Mat}(n,n,\mathbb{F})$ be unitary. Then $\abs{\det U}=1$. More generally, if $S\in\LL(\mathcal{V})$ is an isometry on a finite-dimensional inner-product space over $\mathbb{F}\in\{\mathbb{R},\mathbb{C}\}$, then $\abs{\det S}=1$.
\begin{proof}
The permutation definition gives $\det(U^*)=\overline{\det U}$. By Theorem \ref{thm:determinant-basic-laws}, taking determinants in $U^*U=I_n$ gives $1=\det(U^*)\det(U)=\overline{\det U}\det U=\abs{\det U}^2$. Hence $\abs{\det U}=1$.\\
For an isometry $S$, choose an orthonormal basis and let $U$ be the matrix of $S$ in that basis. Then $U$ is unitary by Theorem \ref{thm:finite-dimensional-isometry-iff-adjoint-product-identity}, and $\det S=\det U$.\\
\hlt{Structural Commentary}: In orthonormal coordinates, isometries are unitary or orthogonal, so the determinant calculation is the matrix one.\\
\hlt{Key Insight}: Determinants turn $U^*U=I$ into $\abs{\det U}^2=1$.
\end{proof}
\end{proposition}

\begin{figure}[H]
  \centering
  \begin{tikzpicture}[
  >=Stealth,
  every node/.style={font=\small},
  box/.style={draw=gray!55, rounded corners=3pt, minimum width=2.5cm, minimum height=0.84cm, align=center, fill=gray!4},
  map/.style={->, line width=0.8pt, blue!65!black}
]
  \node[box] (beta) at (0,0) {$[T]_{\beta\leftarrow\beta}=A$};
  \node[box] (gamma) at (4.2,0) {$[T]_{\gamma\leftarrow\gamma}=PAP^{-1}$};

  \draw[map] (beta) -- (gamma)
    node[midway, above] {$P=C_{\gamma\leftarrow\beta}$};
\end{tikzpicture}

  \caption{Basis changes conjugate matrices; trace and determinant are unchanged.}
  \label{fig:operator-trace-determinant-basis-invariance}
\end{figure}

\begin{proposition} \label{prop:operator-trace-determinant-characteristic-coefficients} \hlt{Trace and Determinant from $q_T$}\\
Let $\mathcal{V}\neq\{\mathbf{0}\}$ be finite-dim over $\mathbb{F}\in\{\mathbb{R},\mathbb{C}\}$, let $T\in\LL(\mathcal{V})$, $q_T(z)=z^n+c_{n-1}z^{n-1}+\cdots+c_1z+c_0$, $n=\dim\mathcal{V}$.\\
Then $\operatorname{tr}T=-c_{n-1}$ and $\det T=(-1)^nc_0$. If $\mathbb{F}=\mathbb{C}$ and $\lambda_1,\ldots,\lambda_n$ list the eigenvalues of $T$ with algebraic multiplicity, then $\operatorname{tr}T=\lambda_1+\cdots+\lambda_n$, $\det T=\lambda_1\cdots\lambda_n$.
\begin{proof}
Choose a basis $\beta$ and set $A=[T]_{\beta\leftarrow\beta}$.\\
Then $q_T=p_A$, and by Definition \ref{def:operator-trace-and-determinant}, $\operatorname{tr}T=\operatorname{tr}A$ and $\det T=\det A$.\\
Theorem \ref{thm:matrix-characteristic-polynomial-basic-coefficients} gives $\operatorname{tr}T=-c_{n-1}$ and $\det T=(-1)^nc_0$.\\
If $\mathbb{F}=\mathbb{C}$, then comparing $q_T(z)=(z-\lambda_1)\cdots(z-\lambda_n)$ with the displayed form of $q_T$ gives the formulas.\\
\hlt{Structural Commentary}: Characteristic polynomial stores trace and determinant in first and last coefficients.\\
\hlt{Key Insight}: Compare coefficients of $\det(zI-T)$.
\end{proof}
\end{proposition}

\begin{proposition} \label{prop:trace-determinant-from-eigenvalues} \hlt{Trace and Determinant from Eigenvalues}\\
Let $\lambda_1,\ldots,\lambda_n$ list the eigenvalues of $A\in\operatorname{Mat}(n,n,\mathbb{C})$ according to $m_A$. Then $\operatorname{tr}A=\sum\limits_{j=1}^{n}\lambda_j$ and $\det A=\prod\limits_{j=1}^{n}\lambda_j$.
\begin{proof}
By Fundamental Theorem of Algebra, $p_A(z)=\prod\limits_{j=1}^{n}(z-\lambda_j) =z^n-\left(\sum\limits_{j=1}^{n}\lambda_j\right)z^{n-1}+\cdots+(-1)^n\prod\limits_{j=1}^{n}\lambda_j$.\\
Theorem \ref{thm:matrix-characteristic-polynomial-basic-coefficients} gives $p_A(z)=z^n-(\operatorname{tr}A)z^{n-1}+\cdots+(-1)^n\det A$. Uniqueness of coefficients gives the identities.\\
\hlt{Structural Commentary}: Coefficients of $p_A$ encode symmetric functions of the eigenvalues.\\
\hlt{Key Insight}: Trace is the first root sum and determinant is the full root product.
\end{proof}
\end{proposition}

\begin{proposition} \label{prop:transpose-conjugate-adjoint-characteristic-polynomial} \hlt{Transpose, Conjugate, and Adjoint Spectra}\\
Let $A\in\operatorname{Mat}(n,n,\mathbb{C})$. Then $p_{A^t}(z)=p_A(z)$, $p_{\overline{A}}(z)=\overline{p_A(\overline{z})}$, and $p_{A^*}(z)=\overline{p_A(\overline{z})}$.\\
Consequently, if $\lambda\in\operatorname{spec}(A)$, then $m_{A^t}(\lambda)=m_A(\lambda)$ and $m_{\overline{A}}(\overline{\lambda})=m_{A^*}(\overline{\lambda})=m_A(\lambda)$.
\begin{proof}
Use determinant identities:
\begin{align}
p_{A^t}(z)&=\det(zI_n-A^t)=\det((zI_n-A)^t)=p_A(z), \nonumber\\
p_{\overline{A}}(z)&=\det(zI_n-\overline{A}) =\det(\overline{\overline{z}I_n-A}) =\overline{\det(\overline{z}I_n-A)} =\overline{p_A(\overline{z})}.
\nonumber
\end{align}
Since $A^*=\overline{A}^{\,t}$, the first identity applied to $\overline{A}$ gives the formula for $A^*$.\\
Multiplicity is preserved because conjugating a polynomial sends a factor $(z-\lambda)^k$ to $(z-\overline{\lambda})^k$.\\
\hlt{Structural Commentary}: Transpose preserves $p_A$; conjugation conjugates its roots.\\
\hlt{Key Insight}: The adjoint combines conjugation with transposition, so its spectrum is conjugated.
\end{proof}
\end{proposition}

\begin{proposition} \label{prop:transpose-preserves-spectrum-and-geometric-multiplicity} \hlt{Transpose Preserves Spectrum and Geometric Multiplicity}\\
Let $A\in\operatorname{Mat}(n,n,\mathbb{F})$ and $\lambda\in\mathbb{F}$. Then $\lambda\in\operatorname{spec}(A)$ $\Longleftrightarrow$ $\lambda\in\operatorname{spec}(A^t)$, and $\dim E_{\lambda}(A)=\dim E_{\lambda}(A^t)$.
\begin{proof}
Since $(A-\lambda I_n)^t=A^t-\lambda I_n$, Theorem \ref{thm:row-rank-equals-column-rank} gives $\operatorname{rank}(A-\lambda I_n)=\operatorname{rank}(A^t-\lambda I_n)$.\\
Matrix rank-nullity then gives $\dim E_{\lambda}(A) =n-\operatorname{rank}(A-\lambda I_n) =n-\operatorname{rank}(A^t-\lambda I_n) =\dim E_{\lambda}(A^t)$.
\begin{enumerate}[label=\roman*.]
\setlength{\itemsep}{0pt}
\item[$(\Rightarrow)$] If $\lambda\in\operatorname{spec}(A)$, then $E_{\lambda}(A)\neq\{\mathbf{0}\}$, so $E_{\lambda}(A^t)\neq\{\mathbf{0}\}$ and $\lambda\in\operatorname{spec}(A^t)$.
\item[$(\Leftarrow)$] The same argument with $A$ and $A^t$ interchanged gives the converse.
\end{enumerate}
\hlt{Structural Commentary}: Transposition preserves rank, so shifted kernels have the same dimension.\\
\hlt{Key Insight}: Spectrum and geometric multiplicity are kernel facts for $A-\lambda I_n$.
\end{proof}
\end{proposition}

\begin{corollary} \label{cor:square-matrix-algebraic-invertibility-dictionary} \hlt{Square Matrix Algebraic Invertibility Dictionary}\\
Let $A=[\mathbf{a}_1\ \cdots\ \mathbf{a}_n]\in\operatorname{Mat}(n,n,\mathbb{F})$. The following are equivalent:
\begin{enumerate}[label=(\roman*)]
\setlength{\itemsep}{0pt}
\item $A$ is invertible;
\item $\mathbf{a}_1,\ldots,\mathbf{a}_n$ form a basis of $\mathbb{F}^n$;
\item $\mathbf{a}_1,\ldots,\mathbf{a}_n$ are linearly independent;
\item $\mathbf{a}_1,\ldots,\mathbf{a}_n$ span $\mathbb{F}^n$;
\item for every $\mathbf{b}\in\mathbb{F}^n$, the system $A\mathbf{x}=\mathbf{b}$ has a unique solution;
\item the homogeneous system $A\mathbf{x}=\mathbf{0}$ has only the zero solution.
\end{enumerate}
\begin{proof}
\begin{enumerate}[label=(\roman*).]
\setlength{\itemsep}{0pt}
\item[$(i)$]$\Rightarrow (ii)$ By Theorem \ref{thm:invertible-matrix-columns-basis}, the columns of an invertible matrix form a basis.
\item[$(ii)$]$\Rightarrow (iii)$ A basis is linearly independent.
\item[$(iii)$]$\Rightarrow (iv)$ In $\mathbb{F}^n$, $n$ linearly independent vectors form a basis, hence span.
\item[$(iv)$]$\Rightarrow (v)$ In $\mathbb{F}^n$, $n$ spanning vectors form a basis, so every $\mathbf{b}$ has unique coordinates $\mathbf{x}$ in $A\mathbf{x}=\mathbf{b}$.
\item[$(v)$]$\Rightarrow (vi)$ Apply (v) to $\mathbf{b}=\mathbf{0}$; since $\mathbf{x}=\mathbf{0}$ is a solution, it is the only one.
\item[$(vi)$]$\Rightarrow (i)$ Condition (vi) says the columns are independent.\\
Hence they form a basis of $\mathbb{F}^n$, and Corollary \ref{cor:basis-columns-form-invertible-matrix} gives invertibility.
\end{enumerate}
\hlt{Structural Commentary}: Square invertibility is column basishood.\\
\hlt{Key Insight}: Invertible columns give unique coordinates.
\end{proof}
\end{corollary}

\begin{corollary} \label{cor:matrix-nonsingularity-rank-determinant-spectrum-dictionary} \hlt{Nonsingularity by Rank, Determinant, and Spectrum}\\
Let $A\in\operatorname{Mat}(n,n,\mathbb{F})$. The following are equivalent:
\begin{enumerate}[label=(\roman*)]
\setlength{\itemsep}{0pt}
\item $A$ is invertible;
\item $\operatorname{rank}(A)=n$;
\item $\det A\neq0$;
\item $0\notin\operatorname{spec}(A)$.
\end{enumerate}
Over $\mathbb{C}$, this is also equivalent to all eigenvalues of $A$ being nonzero.
\begin{proof}
\begin{enumerate}[label=(\roman*).]
\setlength{\itemsep}{0pt}
\item[$(i)$]$\Rightarrow (ii)$ This is Corollary \ref{cor:square-matrix-invertible-iff-full-rank}.
\item[$(ii)$]$\Rightarrow (iii)$ Full rank gives invertibility by Corollary \ref{cor:square-matrix-invertible-iff-full-rank}, hence $\det A\neq0$ by Theorem \ref{thm:determinant-basic-laws}.
\item[$(iii)$]$\Rightarrow (iv)$ If $\det A\neq0$, then $A$ is invertible by Theorem \ref{thm:determinant-basic-laws}; hence $0\notin\operatorname{spec}(A)$ by Corollary \ref{cor:matrix-invertible-iff-zero-not-spectrum}.
\item[$(iv)$]$\Rightarrow (i)$ This is Corollary \ref{cor:matrix-invertible-iff-zero-not-spectrum}.
\end{enumerate}
Over $\mathbb{C}$, Theorem \ref{thm:matrix-characteristic-polynomial-basic-coefficients} identifies $\lambda_1,\ldots,\lambda_n$ as the roots of $p_A$, so (iv) is equivalent to $\lambda_j\neq0$ for all $j$.\\
\hlt{Structural Commentary}: Rank, determinant, and spectrum detect singularity.\\
\hlt{Key Insight}: Singularity can be seen as rank loss, zero determinant, or the zero eigenvalue.
\end{proof}
\end{corollary}

\begin{method} \label{meth:trace-determinant-invariant-computation} \hlt{Trace--Determinant Invariant Method}\\
To get eigenvalue data without solving for eigenvalues:
\begin{enumerate}[label=(\roman*)]
\setlength{\itemsep}{0pt}
\item choose a basis in which $[T]_{\beta\leftarrow\beta}$ is easiest to compute;
\item compute $\operatorname{tr}T$ and $\det T$ from this matrix;
\item if the eigenvalues are $\lambda_1,\ldots,\lambda_n$, conclude $\lambda_1+\cdots+\lambda_n=\operatorname{tr}T$, and $\lambda_1\cdots\lambda_n=\det T$;
\item for equations $ST-TS=A$, first test $\operatorname{tr}A=0$; otherwise no solution exists.
\end{enumerate}
\end{method}

\begin{theorem} \label{thm:rank-product-and-concatenation-bounds} \hlt{Rank Bounds for Products and Concatenations}\\
Let $A\in\operatorname{Mat}(m,k,\mathbb{F})$, $B\in\operatorname{Mat}(k,n,\mathbb{F})$, and $C\in\operatorname{Mat}(m,p,\mathbb{F})$. Then $\operatorname{rank}(AB)\leq \min(\operatorname{rank}(A),\operatorname{rank}(B))$ and $\operatorname{rank}(AB)=\operatorname{rank}(A)$ if and only if $\operatorname{col}(AB)=\operatorname{col}(A)$.\\
Also, $\max(\operatorname{rank}(A),\operatorname{rank}(C))\leq\operatorname{rank}\!\begin{bmatrix}A & C\end{bmatrix}$ and $\operatorname{rank}\!\begin{bmatrix}A & C\end{bmatrix}=\operatorname{rank}(A)$ if and only if $\operatorname{col}(A)=\operatorname{col}(A)+\operatorname{col}(C)$.
\begin{proof}
Each column of $AB$ is a linear combination of the columns of $A$, so $\operatorname{col}(AB)\subseteq\operatorname{col}(A)$. Therefore $\operatorname{rank}(AB)\leq\operatorname{rank}(A)$; since both spaces are finite-dimensional, equality holds exactly when $\operatorname{col}(AB)=\operatorname{col}(A)$. Applying this same argument to $(AB)^t=B^tA^t$ and using Theorem \ref{thm:row-rank-equals-column-rank} gives $\operatorname{rank}(AB)\leq\operatorname{rank}(B)$.\\
For the block matrix, its columns are precisely the columns of $A$ and $C$, hence $\operatorname{col}\!\begin{bmatrix}A & C\end{bmatrix} =\operatorname{col}(A)+\operatorname{col}(C)$.\\
This space contains both $\operatorname{col}(A)$ and $\operatorname{col}(C)$, giving the lower bound. Equality with $\operatorname{rank}(A)$ occurs exactly when the containing subspace has the same dimension as $\operatorname{col}(A)$, equivalently when $\operatorname{col}(A)+\operatorname{col}(C)=\operatorname{col}(A)$.\\
\hlt{Structural Commentary}: Products can only reduce accessible column directions, while horizontal concatenation can only add column directions.\\
\hlt{Key Insight}: Rank tracks whether a matrix operation preserves, loses, or adds column-space dimension.
\end{proof}
\end{theorem}

\begin{theorem} \label{thm:sylvester-rank-inequality} \hlt{Sylvester Rank Inequality}\\
Let $A\in\operatorname{Mat}(m,k,\mathbb{F})$ and $B\in\operatorname{Mat}(k,n,\mathbb{F})$. Then $\operatorname{rank}(A)+\operatorname{rank}(B)-k\leq\operatorname{rank}(AB)$.
\begin{proof}
Let $R=\operatorname{col}(B)\subseteq\mathbb{F}^k$. The columns of $AB$ are $A$ applied to the columns of $B$, so $\operatorname{col}(AB)=A(R)$.\\
Apply rank-nullity to the restriction $T_A|_R:R\to\mathbb{F}^m$. Its range is $A(R)=\operatorname{col}(AB)$ and its kernel is $R\cap\ker(A)$.\\
Hence $\operatorname{rank}(AB) =\dim(A(R)) =\dim(R)-\dim(R\cap\ker(A)) \geq \operatorname{rank}(B)-\dim(\ker(A))$.\\
Since $A$ has $k$ columns, Corollary \ref{cor:matrix-rank-nullity} gives $\dim(\ker(A))=k-\operatorname{rank}(A)$. Substitution gives desired inequality.\\
\hlt{Structural Commentary}: Multiplication by $A$ can lose only the part of $\operatorname{col}(B)$ lying in $\ker(A)$. Sylvester's inequality is this loss estimate in dimensions.\\
\hlt{Key Insight}: The shared middle dimension $k$ bounds how much rank can disappear in the product.
\end{proof}
\end{theorem}

\begin{theorem} \label{thm:frobenius-rank-inequality} \hlt{Frobenius Rank Inequality}\\
Let $A\in\operatorname{Mat}(m,k,\mathbb{F})$, $B\in\operatorname{Mat}(k,p,\mathbb{F})$, and $C\in\operatorname{Mat}(p,n,\mathbb{F})$.\\
Then $\operatorname{rank}(AB)+\operatorname{rank}(BC) \leq \operatorname{rank}(B)+\operatorname{rank}(ABC)$.
\begin{proof}
Let $R=\operatorname{col}(B)$ and $S=\operatorname{col}(BC)$. Since every column of $BC$ is a linear combination of columns of $B$, we have $S\subseteq R$. The induced map $T_A$ maps $R$ onto $\operatorname{col}(AB)$ and maps $S$ onto $\operatorname{col}(ABC)$.\\
Rank-nullity on the restrictions $T_A|_R$ and $T_A|_S$ gives
\begin{align}
\operatorname{rank}(AB)&=\dim R-\dim(R\cap\ker(A)), \quad \operatorname{rank}(ABC)=\dim S-\dim(S\cap\ker(A)). \nonumber
\end{align}
Since $S\cap\ker(A)\subseteq R\cap\ker(A)$, we have $\dim(S\cap\ker(A))\leq\dim(R\cap\ker(A))$.\\
Thus $\dim R-\dim(R\cap\ker(A))+\dim S \leq \dim R+\dim S-\dim(S\cap\ker(A))$.\\
Substituting $\dim R=\operatorname{rank}(B)$ and $\dim S=\operatorname{rank}(BC)$ gives the stated inequality.\\
\hlt{Structural Commentary}: Frobenius tracks $A$ on nested column spaces.\\
\hlt{Key Insight}: The smaller column space cannot lose more dimensions to $\ker(A)$ than the larger one.
\end{proof}
\end{theorem}

\begin{corollary} \label{cor:frobenius-rank-inequality-special-cases} \hlt{Rank Inequalities from Frobenius}\\
The following are special cases of Theorem \ref{thm:frobenius-rank-inequality}:
\begin{enumerate}[label=(\roman*)]
\setlength{\itemsep}{0pt}
\item
$\operatorname{rank}(AB)\leq\min(\operatorname{rank}(A),\operatorname{rank}(B))$;
\item
$\operatorname{rank}(A)+\operatorname{rank}(B)-k\leq\operatorname{rank}(AB)$.
\end{enumerate}
\begin{proof}
\begin{enumerate}[label=(\roman*)]
\setlength{\itemsep}{0pt}
\item
Apply Frobenius to conformable factors $X,Y,Z$. Taking $Z=0$ gives $\operatorname{rank}(XY)\leq\operatorname{rank}(Y)$, because $\operatorname{rank}(Y0)=\operatorname{rank}(XY0)=0$. Taking $X=0$ gives $\operatorname{rank}(YZ)\leq\operatorname{rank}(Y)$.\\
Renaming the factors gives $\operatorname{rank}(AB)\leq\operatorname{rank}(B)$ and $\operatorname{rank}(AB)\leq\operatorname{rank}(A)$.
\item
Taking the middle matrix to be $I_k$ gives $\operatorname{rank}(A)+\operatorname{rank}(B)\leq k+\operatorname{rank}(AB)$, which is Sylvester's inequality.
\end{enumerate}
\hlt{Structural Commentary}: Frobenius is the three-factor parent inequality. The usual upper and lower product-rank bounds appear by making one factor zero or identity.\\
\hlt{Key Insight}: Zero factors isolate upper bounds; an identity middle factor yields the Sylvester lower bound.
\end{proof}
\end{corollary}

\begin{proposition} \label{prop:sylvester-law-of-nullity} \hlt{Sylvester's Law of Nullity}\\
For a matrix $X$, let $\operatorname{nullity}(X)=\dim(\ker(X))$; this is the dimension of its nullspace.
\begin{enumerate}[label=(\roman*)]
\setlength{\itemsep}{0pt}
\item If $A\in\operatorname{Mat}(m,k,\mathbb{F})$ and $B\in\operatorname{Mat}(k,n,\mathbb{F})$, then $\operatorname{nullity}(AB)\leq\operatorname{nullity}(A)+\operatorname{nullity}(B)$.
\item
If $A,B\in\operatorname{Mat}(n,n,\mathbb{F})$, then $\max(\operatorname{nullity}(A),\operatorname{nullity}(B))\leq\operatorname{nullity}(AB)$.
\end{enumerate}
\begin{proof}
\begin{enumerate}[label=(\roman*)]
\setlength{\itemsep}{0pt}
\item
By rank-nullity and Theorem \ref{thm:sylvester-rank-inequality},
\begin{align*}
\operatorname{nullity}(AB) &=n-\operatorname{rank}(AB) \leq n-\operatorname{rank}(A)-\operatorname{rank}(B)+k \\
&=(k-\operatorname{rank}(A))+(n-\operatorname{rank}(B)) =\operatorname{nullity}(A)+\operatorname{nullity}(B).
\end{align*}
\item For square $A,B$, $\operatorname{rank}(AB)\leq\operatorname{rank}(A)$, $\operatorname{rank}(AB)\leq\operatorname{rank}(B)$. Subtract from $n$ gives desired lower bounds.
\end{enumerate}
\hlt{Structural Commentary}: Nullity is rank loss from the domain.\\
\hlt{Key Insight}: Use rank-nullity to translate rank bounds into kernel bounds.
\end{proof}
\end{proposition}

\begin{theorem} \label{thm:rank-triangle-inequality} \hlt{Rank Triangle Inequality}\\
If $A,B\in\operatorname{Mat}(m,n,\mathbb{F})$, then $\abs{\operatorname{rank}(A)-\operatorname{rank}(B)} \leq \operatorname{rank}(A+B) \leq \operatorname{rank}(A)+\operatorname{rank}(B)$.
\begin{proof}
Since $\operatorname{col}(A+B)\subseteq\operatorname{col}(A)+\operatorname{col}(B)$, Theorem \ref{thm:dimension-formula-for-sum} gives $\operatorname{rank}(A+B)\leq\dim(\operatorname{col}(A)+\operatorname{col}(B))\leq\operatorname{rank}(A)+\operatorname{rank}(B)$.
Applying this upper bound to $A=(A+B)-B$ gives $\operatorname{rank}(A)\leq\operatorname{rank}(A+B)+\operatorname{rank}(B)$, so $\operatorname{rank}(A)-\operatorname{rank}(B)\leq\operatorname{rank}(A+B)$. Applying it to $B=(A+B)-A$ gives $\operatorname{rank}(B)-\operatorname{rank}(A)\leq\operatorname{rank}(A+B)$.\\
Together these are exactly $\abs{\operatorname{rank}(A)-\operatorname{rank}(B)}\leq\operatorname{rank}(A+B)$.\\
\hlt{Structural Commentary}: The rank of a sum is controlled by the sum of the two column spaces. The reverse bounds come by solving for either summand.\\
\hlt{Key Insight}: Rank behaves like a dimension-valued size under addition.
\end{proof}
\end{theorem}

